% =============================================================================
% APPENDIX: REPRODUCIBILITY
% =============================================================================
\appendix
\section{Reproducibility Details}
\label{app:repro}
This appendix specifies the exact splits, seeds, model versions, hyperparameters, and code needed to reproduce every reported number, and states the provenance of each value in Table~\ref{tab:comparative}.

\subsection{Datasets and Splits}
\textbf{PlantVillage (controlled, Table~\ref{tab:comparative}).} We use the 11-class novel subset (Apple, Grape, Potato; 9{,}385 images) with frozen embeddings cached once. \emph{11-class regime:} 30 Monte-Carlo episodes; episode $e$ uses \texttt{random.Random(42+17e)}; for each class the image indices are shuffled, the first 5 form the support and all remaining images form the query (9{,}330 queries/episode; support $\cap$ query $=\emptyset$). \emph{5-way $K$-shot regime:} 50 episodes per $K\in\{1,5,10,15\}$; a single \texttt{random.Random(42)} stream samples 5 of the 11 classes per episode (\texttt{rng.sample(range(11),5)}), shuffles each class, takes the first $K$ as support and the rest as query. All methods share identical support/query indices and cached backbone outputs within each episode. Zero-shot BioCLIP and Tip-Adapter use the BioCLIP branches; the centroid and Fisher classifiers use joint BioCLIP+DINOv2 features. Comparisons within the latter group isolate the adaptation rule.

\textbf{PlantDoc (field domain, labeled target-domain support).} 5-way 5-shot: 50 episodes, \texttt{random.Random(42)}, \texttt{rng.sample(range(27),5)}; 5 support leaves/class drawn from the training split, all test images of the sampled classes as queries. 27-way: 135 support leaves (5/class, 8-view augmented to 1{,}080 feature vectors) and 252 field queries. The full PlantDoc dataset is obtained from its original source (not redistributed here); the controlled PlantVillage re-runs require only the cached embeddings and the provided script.

\subsection{Full Shot Ladder}
\label{app:shot_ladder}
Table~\ref{tab:shot_ladder} retains the 10-shot and 15-shot results omitted from the main comparison, alongside the 1-shot and 5-shot entries for reference. Controlled results use 50 episodes per shot count and the splits specified above. Zero-shot BioCLIP uses no support labels; its small variation across shot counts reflects changes in the query set. Published values remain separate because the underlying encoders and class pools differ. No numerical values have been recomputed for this table.

\begin{table*}[htbp]
\caption{Full PlantVillage 5-way shot ladder (accuracy, \%). Controlled results are mean $\pm$ std over 50 episodes per setting. Published results are retained as reported and provide context only. Dashes indicate unreported settings.}
\label{tab:shot_ladder}
\centering
\footnotesize
\setlength{\tabcolsep}{4pt}
\begin{tabularx}{\textwidth}{@{}X c c c c@{}}
\toprule
\textbf{Method} & \textbf{1-shot} & \textbf{5-shot} & \textbf{10-shot} & \textbf{15-shot} \\
\midrule
\multicolumn{5}{@{}l}{\textit{Controlled evaluations, 11-class pool}} \\
Zero-shot BioCLIP & 38.89$\pm$12.78 & 38.86$\pm$12.79 & 38.83$\pm$12.81 & 38.79$\pm$12.83 \\
Tip-Adapter & 45.15$\pm$13.42 & 68.63$\pm$12.80 & 86.01$\pm$7.49 & 93.23$\pm$4.00 \\
Prototypical Networks / SimpleShot & 88.44$\pm$9.13 & 97.45$\pm$2.03 & 98.35$\pm$1.45 & 98.57$\pm$1.27 \\
BioDINO-Edge (inductive) & 88.45$\pm$9.12 & 97.96$\pm$1.85 & 99.10$\pm$0.87 & 99.41$\pm$0.60 \\
BioDINO-Edge (transductive) & 93.42$\pm$6.83 & 98.77$\pm$1.68 & 99.41$\pm$0.60 & 99.63$\pm$0.39 \\
\midrule
\multicolumn{5}{@{}l}{\textit{Published results (context only)}} \\
MAFDE-DN4 \cite{zhao2024mafde} & 57.50 & 81.41 & -- & -- \\
SC-FSL \cite{mu2024few} & 78.55$\pm$0.81 & 92.90$\pm$0.47 & -- & -- \\
AgriFewNet \cite{babu2025agrifewnet} & 87.30 & 94.80 & 97.10 & -- \\
VSF \cite{zhou2026vsf} & 79.01 & 90.51 & -- & -- \\
\bottomrule
\end{tabularx}
\end{table*}

\subsection{Models and Hyperparameters}
Frozen backbones (no fine-tuning): DINOv2 ViT-S/14 (\texttt{timm} \path{vit_small_patch14_dinov2}, 384-D) and BioCLIP ViT-B/16 (\texttt{open\_clip} \path{hf-hub:imageomics/bioclip}, 512-D); joint feature $z=[\,\text{bio}\,\|\,\text{dino}\,]\in\mathbb{R}^{896}$. Fixed constants: logit scale $\gamma=20.0$, Fisher blend $w_{\text{lda}}=0.75$, routing gate $\tau=0.35$; transductive support weight $3\times$, max anchors 20/class, confidence threshold 0.70 on PlantVillage and the buffer settings stated in Section~\ref{sec:metrics} for PlantDoc; Tip-Adapter $\alpha=1.0,\beta=5.5$. Fisher-LDA uses Ledoit--Wolf shrinkage (\texttt{sklearn LDA solver=lsqr, shrinkage=auto}). The estimator applies Ledoit--Wolf per class in standardized feature space and pools the shrunk within-class covariances with class priors; the augmentation sensitivity study (Section~\ref{sec:aug_sensitivity}) reports the mean per-class intensity. At 5-way 1-shot the raw support has $N=C=5$ and the within-class scatter vanishes; we resolve the singularity with a data-driven ridge floor (mean support feature variance) added to the Ledoit--Wolf covariance, which reduces Fisher to regularized nearest-centroid (the correct degenerate limit). For $N>C$ the canonical sklearn estimator is used unchanged.

\textbf{Why two solvers appear in the paper.} Section~\ref{sec:shrinkage} writes the estimator analytically as a pooled scatter matrix with a Cholesky solve, whereas the Python reference uses \texttt{sklearn} with \texttt{solver=lsqr} and an automatically selected shrinkage, and the Android client uses a dense Cholesky factorization in EJML. These are three routes to the same regularized discriminant: \texttt{lsqr} solves the same normal equations by conjugate gradients rather than by explicit factorization, so with a converged tolerance it returns the pooled Ledoit--Wolf solution up to floating point. The on-device timings in Table~\ref{tab:adaptation_breakdown} (332.00\,ms Cholesky, 51.31\,ms projection solve) therefore reflect the deterministic dense path used on the phone, and agreement between the two is what the FP32/INT8 comparison in Table~\ref{tab:quantization} verifies: the same split, the same constants, and a difference of at most 0.40\,pp on either dataset.

\subsection{Environment and Code}
Reference environment: Python 3.12.1, torch 2.6.0 (Compute Unified Device Architecture (CUDA) 12.4 build; all reported experiments execute on CPU), scikit-learn 1.9.1, numpy 2.5.4, scipy 1.18.1, timm 1.0.30, open\_clip 3.3.0. The controlled comparison is reproduced by \path{run_controlled_baselines.py} in \path{experiments/baseline_comparison/} (output \path{controlled_baseline_results.json}); the canonical PlantVillage reference run by \path{benchmark_engine.py} (\path{pipeline/}) and \path{run_transductive_validation.py} (\path{experiments/transductive_validation/}); the augmentation sensitivity study by \path{run_augmentation_sensitivity.py} (\path{experiments/augmentation_sensitivity/}, output \path{results/augmentation_sensitivity.json}; augmentation RNG re-seeded per episode and pass from the episode seed with fixed offsets). Code and pinned release: \url{https://github.com/reyhane-code/BioDINO-Edge-android}.

\subsection{Provenance and Numerical Reproducibility}
Table~\ref{tab:provenance} states the origin of each reported value (``Re-run'': computed by us with the provided script on the stated protocol; ``As reported'': quoted from the original publication, not re-run). The controlled script reproduces the canonical inductive accuracy exactly (95.08\%) and the transductive accuracy within 0.04\,pp (97.37\% re-run vs.\ 97.41\% reference); the residual is floating-point sensitivity of the iterative least-squares (\texttt{lsqr}) LDA solve across library versions, far below the episode standard deviation (0.71\,pp). The centroid/cosine baselines involve no iterative solve and reproduce exactly. We report the canonical reference values for the 11-class regime to remain consistent with the paired-significance and sensitivity analyses (Sections~\ref{sec:ablation}--\ref{sec:aug_sensitivity}).

\begin{table}[htbp]
\caption{Provenance of reported values.}
\label{tab:provenance}
\centering

\footnotesize
\begin{tabularx}{\columnwidth}{l X}
\toprule
\textbf{Value} & \textbf{Provenance} \\
\midrule
Table~\ref{tab:comparative} baselines & Re-run, shared splits/seeds; feature branches follow each method \\
Table~\ref{tab:comparative} BioDINO 11-class & Canonical reference run (reproduced $\le$0.04\,pp) \\
Table~\ref{tab:comparative} BioDINO 5-way & Controlled re-run (this work) \\
Table~\ref{tab:shot_ladder} literature rows & As reported, not re-run \\
Table~\ref{tab:plantdoc_context} & Our PlantDoc field-domain run (seeded 8-view ladder, protocol above) \\
Table~\ref{tab:aug_sensitivity} sensitivity & Re-run, this work (same 30-episode protocol) \\
\bottomrule
\end{tabularx}
\end{table}

